% ==============================================================================
% SETUP.TEX - CONFIGURACIÓN GLOBAL DEL PROYECTO TC2
% ==============================================================================
\ifdefined\SETUPCONFIGURED
\else
\newcommand{\SETUPCONFIGURED}{}

% Este archivo centraliza todos los paquetes y estilos.
% Debe incluirse en el preámbulo del libro principal y de los archivos standalone.

% --- 1. CODIFICACIÓN E IDIOMA ---
\usepackage[utf8]{inputenc}
\usepackage[T1]{fontenc}
\usepackage[spanish, es-tabla]{babel}  
\usepackage{csquotes} 

\usepackage{import}
\usepackage[mode=image, subpreambles=false]{standalone}

% --- 2. MATEMÁTICAS Y CIENCIAS ---
\usepackage{amsmath, amssymb, amsfonts, mathtools}

% Operadores en notación español/europea (seno, coseno y seno hiperbólicos)
\DeclareMathOperator{\sen}{sen}
\DeclareMathOperator{\ch}{ch}
\DeclareMathOperator{\sh}{sh}

% --- 3. DISEÑO Y GEOMETRÍA --- 
\usepackage[left=2.5cm,right=2.5cm,top=3cm,bottom=3cm]{geometry}
\usepackage{parskip} 
\usepackage{float}   

% --- 4. GRÁFICOS Y TABLAS ---
\usepackage{graphicx}
\usepackage{booktabs} 
\usepackage{tikz}
\usetikzlibrary{arrows.meta, calc, angles, quotes, babel, shapes.geometric}
\usepackage{pgfplots}
\usepgfplotslibrary{groupplots}
\usepgfplotslibrary{external}
\usepgfplotslibrary{patchplots}
\pgfplotsset{compat=1.18}
\usepackage[siunitx, RPvoltages]{circuitikz}

% --- 5. COLORES Y CAJAS ---
\usepackage[table]{xcolor}
\usepackage[theorems, most]{tcolorbox}

\definecolor{utnblue}{RGB}{0, 51, 102}
\definecolor{codegreen}{rgb}{0,0.6,0}
\definecolor{codegray}{rgb}{0.5,0.5,0.5}
\definecolor{codepurple}{rgb}{0.58,0,0.82}
\definecolor{backcolour}{rgb}{0.96,0.96,0.96}

% --- 6. ENTORNOS PERSONALIZADOS (CAJAS) ---
\newtcolorbox{pedagogico}{
    colback=codegreen!5!white,
    colframe=codegreen,
    title=\textbf{Enfoque Pedagógico},
    fonttitle=\bfseries,
    sharp corners,
    boxrule=0.5mm
}

\newtcolorbox{ingenieria}{
    colback=utnblue!5!white,
    colframe=utnblue,
    title=\textbf{Utilidad Ingenieril},
    fonttitle=\bfseries,
    sharp corners,
    boxrule=0.5mm
}

\newtcolorbox{nota}{
    colback=codegray!5!white,
    colframe=codegray,
    title=\textbf{Nota},
    fonttitle=\bfseries,
    sharp corners,
    boxrule=0.5mm
}

% --- 7. CÓDIGO FUENTE (LISTINGS) ---
\usepackage{listings}
\lstdefinestyle{miestilo}{
    backgroundcolor=\color{backcolour},   
    commentstyle=\color{codegreen},
    keywordstyle=\color{magenta},
    numberstyle=\tiny\color{codegray},
    stringstyle=\color{codepurple},
    basicstyle=\ttfamily\footnotesize,
    breakatwhitespace=false,         
    breaklines=true,                 
    captionpos=b,                    
    keepspaces=true,                 
    numbers=left,                    
    numbersep=5pt,                  
    showspaces=false,                
    showstringspaces=false,
    showtabs=false,                  
    tabsize=2,
    frame=single,                    
    rulecolor=\color{codegray}
}
\lstset{style=miestilo}
\lstset{literate={ñ}{{\~n}}1 {á}{{\'a}}1 {é}{{\'e}}1 {í}{{\'i}}1 {ó}{{\'o}}1 {ú}{{\'u}}1}

% --- 8. HIPERVÍNCULOS ---
\usepackage[colorlinks=true, linkcolor=utnblue, urlcolor=utnblue, citecolor=utnblue]{hyperref}

% --- 9. PLACEHOLDER PARA FIGURAS PENDIENTES ---
% Marca el lugar y la descripción de una figura aún no dibujada (p. ej. Cap. 5:
% circuitos y gráficas diferidos). Greppable por \figpendiente para el pase posterior.
\newcommand{\figpendiente}[1]{%
  \begin{center}%
  \setlength{\fboxsep}{6pt}%
  \fbox{\parbox{0.85\linewidth}{\centering\small\itshape
  [\textbf{Figura pendiente}]\\[2pt] #1}}%
  \end{center}}

\fi
